% =====================================================================
%  Thème 3 — Relation fondamentale & formules d'addition — NIVEAU FACILE
%  Exercices
% =====================================================================
\documentclass[11pt,a4paper]{article}
\usepackage[utf8]{inputenc}
\usepackage[T1]{fontenc}
\usepackage{lmodern}
\usepackage[french]{babel}
\usepackage{amsmath,amssymb,mathtools}
\usepackage{icomma}
\usepackage{xcolor}
\usepackage{colortbl}
\usepackage{tikz}
\usepackage[most]{tcolorbox}
\usepackage{enumitem}
\usepackage{array}
\usepackage[a4paper,margin=2.2cm]{geometry}
\usetikzlibrary{arrows.meta,calc}

\definecolor{primary}{RGB}{214,51,108}
\definecolor{primarydark}{RGB}{140,20,64}

\DeclareMathOperator{\tg}{tg}
\DeclareMathOperator{\cotg}{cotg}
\newcommand{\degr}{^{\circ}}
\newcommand{\Z}{\mathbb{Z}}

\newtcolorbox{consigne}{enhanced,breakable,colback=primary!5,colframe=primary,
  boxrule=0.8pt,arc=2pt,left=3mm,right=3mm,top=2mm,bottom=2mm}

\newcounter{exo}
\newcommand{\exo}[1]{\medskip\refstepcounter{exo}%
  \noindent{\large\bfseries\color{primarydark}Exercice \theexo}%
  \ \textcolor{primary}{\textbullet}\ \textbf{#1}\par\nopagebreak\smallskip}

\setlength{\parindent}{0pt}
\pagestyle{empty}

\begin{document}

\begin{center}
{\LARGE\bfseries\color{primary}Thème 3 — Relation fondamentale \& formules d'addition}\\[2pt]
{\color{primarydark}\textsc{Niveau Facile — Exercices}}
\end{center}
\hrule height 1pt
\medskip

\begin{consigne}
\textbf{Objectif :} appliquer directement la relation fondamentale, les trois écritures de $\cos2a$ et
une formule d'addition sur des angles remarquables. \textbf{Sans calculatrice}, valeurs
\emph{exactes}. Quand un signe $\pm$ apparaît, le quadrant est toujours précisé.
\end{consigne}

\exo{La relation fondamentale de tête}
On donne $\cos\alpha$ ; calcule $\sin^2\alpha$ à l'aide de $\cos^2\alpha+\sin^2\alpha=1$.
\[
\textbf{(a)}\ \cos\alpha=\frac{3}{5} \qquad
\textbf{(b)}\ \cos\alpha=\frac{\sqrt3}{2} \qquad
\textbf{(c)}\ \cos\alpha=\frac{1}{3} \qquad
\textbf{(d)}\ \cos\alpha=0
\]

\exo{Le signe décidé par le quadrant}
Détermine $\cos\alpha$ (valeur exacte, signe compris).
\[
\textbf{(a)}\ \sin\alpha=\frac{4}{5},\ \text{quadrant I} \qquad
\textbf{(b)}\ \sin\alpha=\frac{1}{2},\ \text{quadrant II} \qquad
\textbf{(c)}\ \sin\alpha=-\frac{\sqrt2}{2},\ \text{quadrant III}
\]

\exo{Reconnaître les trois visages de $\cos 2a$}
Complète chaque égalité par la bonne expression ($\cos^2a-\sin^2a$, \ $2\cos^2a-1$ \ ou \ $1-2\sin^2a$).
\[
\textbf{(a)}\ \cos2a=\ldots\ \text{(avec seulement }\cos a) \qquad
\textbf{(b)}\ \cos2a=\ldots\ \text{(avec seulement }\sin a)
\]
\[
\textbf{(c)}\ \text{Écris }\sin2a\text{ en fonction de }\sin a\text{ et }\cos a.
\]

\exo{Une addition sur des angles connus}
En utilisant $\cos(a+b)=\cos a\cos b-\sin a\sin b$, calcule (les deux angles sont remarquables).
\[
\textbf{(a)}\ \cos\!\Big(\frac{\pi}{3}+\frac{\pi}{6}\Big) \qquad
\textbf{(b)}\ \sin\!\Big(\frac{\pi}{4}+\frac{\pi}{4}\Big) \qquad
\textbf{(c)}\ \cos\!\Big(\frac{\pi}{2}+\frac{\pi}{3}\Big)
\]

\exo{Linéariser}
À l'aide de $\cos^2x=\dfrac{1+\cos2x}{2}$ et $\sin^2x=\dfrac{1-\cos2x}{2}$, exprime en fonction de
$\cos2x$ (sans carré).
\[
\textbf{(a)}\ \cos^2x \qquad
\textbf{(b)}\ \sin^2x \qquad
\textbf{(c)}\ \cos^2x-\sin^2x
\]

\end{document}
